\section{治理建议}
问题一中，四大家鱼恢复上升，水草、浮游动物和底栖饵料却可能下降；问题二中，江豚与长江鲟对通道和放流的响应又不一样。治理重点不能只停留在捕捞约束本身。后续监测最好同时列出四大家鱼总量、底层资源指标和食压指数；鱼道修复、产卵场连通及放流节律也不宜按同一办法处理。

更迫切的，是把污染治理和入侵控制往前提。问题五已经说明，污染情景的部分食源过程指标短时上升，珍稀物种终态和综合得分却在下降，所以局部过程变好并不等于系统整体好转。对入侵物种，可把鄱阳湖等支流湖泊作为高风险前沿，优先布设早期监测、拦截和清除，避免其沿食物网空档扩散。禁渔只是减轻了捕捞压力；如果水质治理、通道修复、放流优化和入侵防控接不上，恢复收益仍留不住。

